\section{Methods and Results（研究方法・結果）}

\subsection{研究対象}

本研究では、ホテルサンバレー那須において実施された地域住民参加型クリスマス聖歌企画を研究対象とする。

本企画は、地域住民による自主的な合唱活動をホテル施設内で行い、さらに観光客自身も歌唱へ参加できる構造を持つものである。参加人数はオルガン演奏者を含め約25名であり、参加者には高齢者も多く含まれていた。

本研究では、高齢者による歌声そのものを、単なる音楽技術としてではなく、長年の人生経験によって形成された表現として捉える。

楽器が長い年月を経て音色を深めていくように、人間の声もまた、

\begin{itemize}
\item 苦労
\item 喜び
\item 地域生活
\item 時代背景
\item 人生経験
\end{itemize}

を積み重ねることで独自の響きを持つ。

そのため本研究では、高齢者の歌声に含まれる「生きてきた時間」を、地域文化資源の一部として観察対象に含める。

また現代社会では、高齢者自身が「自分はもう役割を終えた存在である」と感じてしまう場面も少なくない。一方で若い世代側には、「その経験こそ必要である」「地域の前に立ってほしい」という感覚も存在している。

本研究では、この両者が接続されることで、

\begin{itemize}
\item 高齢者の役割再形成
\item 経験知の地域資源化
\item 長寿命社会における社会参加
\end{itemize}

などがどのように形成されるかについても観察対象とする。

さらに本研究では、「観光」の定義そのものについても再考を行う。

従来の観光は、

\begin{itemize}
\item 景色を見る
\item 宿泊する
\item 食事を楽しむ
\end{itemize}

など、消費行動を中心として語られることが多い。

しかし本企画では、

\begin{itemize}
\item 地域住民と共に歌う
\item 他の観光客と空間を共有する
\item 同じ時間を過ごす
\end{itemize}

という体験を通じて、「この場所に愛着を持つ」という感覚がどのように形成されるかを観察対象とする。

特に本研究では、

\begin{itemize}
\item 観光客の笑顔
\item 家族間の雰囲気
\item 自然な会話
\item 空間内での一体感
\end{itemize}

など、数値化しにくい感情的・空間的反応についても研究対象として扱う。

また本研究では、歌詞カード・曲解説・感謝メッセージなどを、単なる補助資料ではなく、人と文化と空間を接続するための「情報設計」として扱う。

ここでいう情報設計とは、単なる情報伝達ではなく、

\begin{itemize}
\item なぜ歌うのか
\item なぜ集まるのか
\item なぜ共に時間を過ごすのか
\end{itemize}

という意味そのものを共有し、人間同士を自然に接続するための社会的設計を含んでいる。

4. Results（結果）


\subsection{地域住民主体による文化形成}

本企画では、外部演者依存ではなく、地域住民自身による文化活動として合唱が形成された。

特に合唱という形式は、多人数参加を必要とするため、通常の商業イベントでは高コスト化しやすい。そのため一般的なホテルイベントや観光施設では、

\begin{itemize}
\item 歌い手が一人または二人
\item 少人数編成
\item 演奏中心構成
\end{itemize}

となる場合が多く、一般の観光客が「合唱」を実際に体験する機会は決して多くない。

しかし本企画では、

\begin{itemize}
\item 趣味活動
\item 地域交流
\item 自主参加
\end{itemize}

を基盤とすることで、多人数による合唱空間が成立していた。

さらに重要なのは、この合唱が単なる音楽演出ではなく、「地域に生きる人々の声」によって構成されていた点である。

特に高齢者を中心とした歌声には、長年の人生経験や地域生活によって形成された深みが含まれており、それが演奏技術とは異なる空間的な温度や安心感を生み出していた。

これは、地域に既に存在していた人的資源や経験知が、観光価値として再接続された事例である。


\subsection{観光客参加による体験変化}

観光客は単なる鑑賞者ではなく、実際に歌唱へ参加した。

その結果、

\begin{itemize}
\item 家族同士で笑顔を共有する様子
\item カップルが共に歌う様子
\item 子どもが自然に参加する様子
\item 観光客同士が互いに微笑み合う様子
\end{itemize}

などが観察された。

特に子連れの観光客では、子どもだけではなく、保護者側も「めちゃくちゃ楽しい」と自然に声を上げる場面が見られた。

これは単なるイベント鑑賞ではなく、「家族全体で同じ体験へ参加する構造」が形成されていたことを示している。

また、観光客が歌唱へ参加することで、

\begin{itemize}
\item 地域住民
\item ホテル関係者
\item 観光客
\end{itemize}

という立場の違いが一時的に薄れ、「共に歌う人」として同じ空間を共有する状態が生まれていた。

特に、「一緒に歌った」という行為そのものが、観光体験を記憶へ残す重要な要素となっていた。


\subsection{文化の体験的理解}

本企画では、単に音楽を演奏するだけではなく、

\begin{itemize}
\item 曲背景
\item 歌の歴史
\item サンタクロースの由来
\end{itemize}

などを短く説明することで、観光客が「なぜこの文化が存在しているのか」を感覚的に理解する構造が形成されていた。

例えば、「さやかに星はきらめく（O Holy Night）」については、当初教会から否定されながらも、旋律の美しさによって世界中へ広がった背景が紹介された。

また、サンタクロースの由来として、困窮した子どもたちを救うために金貨を靴下へ入れた逸話なども紹介された。

観光客からは「なるほど」という反応や笑顔も見られ、これらは単なる知識説明ではなく、「文化へ参加する理由」を共有する役割を持っていた。

さらに本企画では、歌詞カード内に掲載された感謝のメッセージを通じ、「訪れる人」「迎える人」「暮らす人」が共に歌い、同じ時間を共有するという意味が明確に伝えられていた。

これは単なるイベント説明ではなく、「なぜこの場を共にするのか」という感覚そのものを共有する情報設計として機能していた。


\subsection{地域参加者側への影響}

本企画では、観光客だけではなく、地域住民側にも感情的変化が観察された。

特に合唱団参加者の一人は、本番後に涙を流していた。

これは単なる演奏達成感ではなく、

\begin{itemize}
\item 地域住民
\item 観光客
\item ホテル
\item 空間
\item 音楽
\end{itemize}

が一つにつながった感覚による感情反応であった可能性がある。

また本企画終了後、ホテル側から「ぜひ毎年続けてほしい」という要望も生まれた。

これは単なる単発イベントとしてではなく、ホテル側にとっても、

\begin{itemize}
\item 地域文化形成
\item 観光価値向上
\item 空間再活用
\item 地域との接続
\end{itemize}

として有効性が認識されたことを示している。

特に重要なのは、本企画が外部大型イベントではなく、地域住民主体によって形成されていた点である。

つまり本研究で観察された価値は、「消費されるイベント」ではなく、「地域と共に育つ文化構造」として認識され始めていた。


\subsection{空間再活用としての成功}

太陽の教会は、元々結婚式場として設計された空間であり、

\begin{itemize}
\item パイプオルガン
\item 教会建築
\item 光演出
\end{itemize}

などを備えていた。

しかし、少子化や結婚式需要の減少によって、以前ほど活用されなくなっていた背景がある。

本企画では、その空間が再び、

\begin{itemize}
\item 歌声
\item 人
\item 文化
\item 感情
\end{itemize}

によって満たされ、本来の空間価値が再接続されていた。

これは単なるイベント利用ではなく、「空間本来の意味」を地域文化によって再活性化した事例である。
